\paragraph{The ordinary real-multiplication restriction is proper.}
There are ordinary genus-two curves over finite fields of characteristic
two whose geometric Jacobians admit no Rosati-invariant action of
$\mathcal O_5$, even after extending constants. An explicit example is
the smooth proper curve over $\mathbf F_4=\mathbf F_2(u)$ given by
\[
 C_0:\qquad y^2+y=x+\frac{u}{x}+\frac{u^2}{x+1},
                         \qquad u^2+u+1=0.
 \tag{DRP.1}
\]
Its Jacobian is absolutely simple and has geometric rational
endomorphism algebra a quartic field whose Rosati-fixed subfield is
$\mathbf Q(\sqrt2)$. Thus the additional restriction in A33 is an
actual proper class of ordinary curves, not merely a formally added
condition. This example makes no claim against point-section
injectivity on any curve.

\begin{proof}
The three poles $0,1,\infty$ in (DRP.1) are simple. The
Artin--Schreier genus and rank formula gives genus two and $2$-rank
two \cite[Lemma~2.6]{PriesZhu2012}. Each pole has one rational point
above it. For $x=u,u^2$, the right side is $u$, whose trace to
$\mathbf F_2$ is one. Therefore $\#C_0(\mathbf F_4)=3$.

For the second count write
$\mathbf F_{16}=\mathbf F_2[z]/(z^4+z+1)$ and $u=z^2+z$.
Encode a field element by its four binary coefficients in increasing
powers of $z$, so $u=6$ and $u^2=7$. Among the fourteen values
$x\in\mathbf F_{16}\setminus\{0,1\}$, the right side of
(DRP.1) has absolute trace zero exactly at
\[
                         x\in\{6,7,8,11,13,15\}.
 \tag{DRP.2}
\]
Each supplies two rational points. Hence
$\#C_0(\mathbf F_{16})=3+2\cdot6=15$.
The accompanying integer-only certificate retains all fourteen
right-side values and traces, and independently enumerates the
solutions of $y^2+y$ to verify the trace counts.

Let $\pi$ be $4$-power Frobenius on $A=J_{C_0}$. The two point
counts and the genus-two Newton identities give its characteristic
polynomial
\[
                    P(X)=X^4-2X^3+X^2-8X+16.
 \tag{DRP.3}
\]
Indeed its cubic coefficient is $3-5=-2$, and its quadratic
coefficient is $((-2)^2-(17-15))/2=1$.
Modulo three, $P$ is $\Phi_5(X)$, irreducible because
$\operatorname{ord}_5(3)=4$. Thus $P$ is irreducible over
$\mathbf Q$, and $K=\mathbf Q(\pi)$ has degree four.

We next exclude root-of-unity ratios between distinct eigenvalues.
The exact polynomial identity
\[
 \begin{aligned}
 \operatorname{Res}_X(P(X),P(tX))
   &=256(t-1)^4 A_0(t)^2 B_0(t),\\
 A_0(t)&=4t^4+7t^3+10t^2+7t+4,\\
 B_0(t)&=16t^4-8t^3+t^2-8t+16
 \end{aligned}
 \tag{DRP.4}
\]
can be checked as the determinant of $P(tM)$, where $M$ is the
four-by-four companion matrix of $P$. The certificate evaluates that
determinant by its complete permutation formula over $\mathbf Z[t]$.
Both $A_0$ and $B_0$ are primitive and reduce modulo three to
$\Phi_5$, so both are irreducible over $\mathbf Q$. Their monic
normalizations have nonintegral coefficients, so their roots are not
algebraic integers and in particular are not roots of unity.
The roots of the resultant are the quotients $\alpha_j/\alpha_i$
of the four roots of $P$. Since $P$ has distinct roots, the four
factors $t-1$ account exactly for the diagonal quotients. No quotient
between distinct roots is a root of unity. Consequently $\pi^m$
has four distinct eigenvalues for every $m\ge1$.

Here the standard integral Tate-module injectivity, rather than mere
faithfulness over $\mathbf Q$, is relevant. Milne's
Theorem~12.5 states that the geometric endomorphism group is a finite
free $\mathbf Z$-module and that its tensor product with
$\mathbf Z_\ell$ injects into the endomorphisms of $T_\ell A$
for $\ell\ne2$ \cite[Theorem~12.5, pp.~21--22]{MilneAV}.
Choose a finite field over which a basis of this endomorphism group
is defined; every basis element is defined over some finite field,
so one common extension exists. Its endomorphisms commute with some
$\pi^m$. On the four-dimensional rational Tate module the
centralizer of $\pi^m$ has dimension four: after extending scalars
it is the diagonal algebra for four distinct eigenvalues. The stated
injectivity therefore gives
\[
 \dim_{\mathbf Q}\operatorname{End}^0_{\overline{\mathbf F}_4}(A)
       \le4.
\]
Since this algebra already contains the quartic field $K$, it equals
$K$. In particular $A$ is absolutely simple: a nontrivial geometric
isogeny decomposition would give a nontrivial idempotent in this
field, by Poincare reducibility
\cite[Proposition~12.1, p.~19]{MilneAV}.

The Rosati involution for the canonical polarization satisfies
$\pi^\dagger=4/\pi$ \cite[Lemma~19.2, p.~41]{MilneAV}.
This is a nontrivial involution on $K$. Put $s=\pi+4/\pi$.
Dividing (DRP.3) by $\pi^2$ gives
\[
                         s^2-2s-7=0.
 \tag{DRP.5}
\]
Its discriminant is $32$, so $\mathbf Q(s)=\mathbf Q(\sqrt2)$.
This quadratic field is fixed by Rosati and is the entire fixed
field, since $[K:\mathbf Q]=4$ and the involution is nontrivial.
A Rosati-invariant embedding of $\mathcal O_5$ would extend to an
embedding of $\mathbf Q(\sqrt5)$ into $\mathbf Q(\sqrt2)$,
which is impossible. The assertion already concerns all geometric
endomorphisms, so no constant extension changes the conclusion.
\end{proof}
